\documentclass[a4paper, 12pt]{amsbook} 
\usepackage{amssymb,amsmath,enumitem,pgfplots,hyperref,setspace,fancyhdr,verbatim,array,multirow, booktabs,array,xtab,pifont,tikzpeople,colortbl,tikz-cd,adjustbox,longtable,tikz-3dplot, tasks}
\usepackage[most]{tcolorbox}
\usetikzlibrary{decorations.markings,arrows,trees,positioning,shapes,fit,calc,decorations.pathreplacing,intersections}
\usepgfplotslibrary{fillbetween}
\setlength{\itemsep}{1in}

\onehalfspacing

\newcommand{\ZeroRoman}[1]{
	\ifcase\value{#1}\relax 0\else
	\Roman{#1}\fi}

\renewcommand{\thechapter}{\ZeroRoman{chapter}}
\renewcommand{\chaptername}{Topic}
\renewcommand{\chaptermark}[1]{\markleft{$\S$\thechapter.\ #1}{}}
\renewcommand{\sectionmark}[1]{\markright{$\S$\thesection.\ #1}{}}
\renewcommand{\thesection}{\thechapter.\arabic{section}}

\newcolumntype{N}{>{\centering\arraybackslash}m{1.8in}}

\setlength{\textwidth}{6.8in}

\calclayout

\definecolor{bananayellow}{rgb}{1.0, 0.88, 0.21}
\definecolor{darkred}{rgb}{0.55, 0.0, 0.0}
\definecolor{darkblue}{rgb}{0.0, 0.0, 0.55}
\definecolor{airforceblue}{rgb}{0.36, 0.54, 0.66}
\definecolor{amethyst}{rgb}{0.6, 0.4, 0.8}
\definecolor{violet(ryb)}{rgb}{0.53, 0.0, 0.69}
\definecolor{arsenic}{rgb}{0.23, 0.27, 0.29}
\definecolor{darkgreen}{rgb}{0.0, 0.2, 0.13}
\definecolor{darkjunglegreen}{rgb}{0.1, 0.14, 0.13}
\definecolor{bazaar}{rgb}{0.6, 0.47, 0.48}
\definecolor{blanchedalmond}{rgb}{1.0, 0.92, 0.8}
\definecolor{babyblue}{rgb}{0.54, 0.81, 0.94}
\definecolor{beaublue}{rgb}{0.74, 0.83, 0.9}
\definecolor{asparagus}{rgb}{0.53, 0.66, 0.42}
\definecolor{darkspringgreen}{rgb}{0.09, 0.45, 0.27}
\definecolor{darktangerine}{rgb}{1.0, 0.66, 0.07}
\definecolor{alizarin}{rgb}{0.82, 0.1, 0.26}

\newtcbtheorem[auto counter]{objectivebox}{Objective}{
	rounded corners,
	breakable,
	colback=blue!5!white,
	colframe=blue!75!black,
	fonttitle=\bfseries,
	description font=\bfseries,
	before skip=0.5cm, % Adjust this to change the space before the box
	after skip=0.5cm,  % Adjust this to change the space after the box
	code={\doublespacing}
}{obj}

\newcommand{\Arrow}[1]{%
	\parbox{#1}{\tikz{\draw[->](0,0)--(#1,0);}}
}

\tcbset{theostyle/.style={
		enhanced,
		sharp corners,
		attach boxed title to top left={
			xshift=-1mm,
			yshift=-4mm,
			yshifttext=-1mm
		},
		top=1.5ex,
		colback=white,
		colframe=blue!75!black,
		fonttitle=\bfseries,
		description font=\normalfont,
		boxed title style={
			sharp corners,
			size=small,
			colback=blue!75!black,
			colframe=blue!75!black,
		},
}}

\newtcbtheorem[no counter]{randombox}{}{
	arc=0pt,
	outer arc=0pt,
	colback=blanchedalmond,
	breakable,
	colframe=bazaar,
	separator sign none,
	before skip=0.5cm, % Adjust this to change the space before the box
	after skip=0.5cm,  % Adjust this to change the space after the box
	code={\doublespacing},
}{rdm}

\fancyhf{}
\fancyhead[RE]{\nouppercase{\leftmark}}
\fancyhead[LO]{\nouppercase{\rightmark}}
\fancyhead[RO]{\thepage}
\fancyhead[LE]{\thepage}

\newtcbtheorem[number within=section]{coloredtheorem}{}%
{colback=beaublue,colframe=babyblue,fonttitle=\bfseries}{th}

\newtcbtheorem[number within=section]{ColoredDefinition}{Definition}{theostyle}{def}

\newcommand{\vasymptote}[2][]{
	\draw [red,densely dashed,#1] ({rel axis cs:0,0} -| {axis cs:#2,0}) -- ({rel axis cs:0,1} -| {axis cs:#2,0});
}

\newcommand{\hasymptote}[2][]{
	\draw [red,densely dashed,#1] ({rel axis cs:0,0} |- {axis cs:0,#2}) -- ({rel axis cs:1,0} |- {axis cs:0,#2});
}

\newcounter{examplecounter}
\counterwithin{examplecounter}{section}

\newtheorem{theorem}{Theorem}%[chapter]
\newtheorem{example}[examplecounter]{Example}
\newtheorem{checkyourunderstanding}{CYU}%[chapter]{CYU}
\newtheorem*{remark}{Remark}

\newcounter{choice}
\renewcommand\thechoice{\textit{\Alph{choice}}}
\newcommand\choicelabel{\thechoice.}

\newenvironment{choices}%
{\list{\choicelabel}%
	{\usecounter{choice}\def\makelabel##1{\hss\llap{##1}}%
		\settowidth{\leftmargin}{W.\hskip\labelsep\hskip 2.5em}%
		\def\choice{%
			\item
		} % task
		\labelwidth\leftmargin\advance\labelwidth-\labelsep
		\topsep=2pt
		\partopsep=2pt
	}%
}%
{\endlist}

\makeatletter
\def\@makechapterhead#1{%
	%% change the value of \topskip from 7.5pc to the desired value
	\global\topskip 0pc\relax
	\begingroup
	\fontsize{\@xivpt}{18}\bfseries\centering
	\ifnum\c@secnumdepth>\m@ne
	\leavevmode \hskip-\leftskip
	\rlap{\vbox to\z@{\vss
			\centerline{\normalsize\mdseries
				\uppercase\@xp{\chaptername}\enspace\thechapter}
			\vskip 3pc}}\hskip\leftskip\fi
	#1\par \endgroup
	\skip@34\p@ \advance\skip@-\normalbaselineskip
	\vskip\skip@ }
\def\@makeschapterhead#1{%
	%% change the value of \topskip from 7.5pc to the desired value
	\global\topskip 0pc\relax
	\begingroup
	\fontsize{\@xivpt}{18}\bfseries\centering
	#1\par \endgroup
	\skip@34\p@ \advance\skip@-\normalbaselineskip
	\vskip\skip@ }
\makeatother
\newcommand{\PreserveBackslash}[1]{\let\temp=\\#1\let\\=\temp}
\newcolumntype{R}[1]{>{\PreserveBackslash\raggedright}p{#1}}
\newcolumntype{L}[1]{>{\PreserveBackslash\raggedleft}p{#1}}
\newcolumntype{M}[1]{>{\centering\arraybackslash}m{#1}}
\pgfdeclarelayer{bg}
\pgfsetlayers{bg,main}

\pagenumbering{gobble}

%\usepackage[a4paper, total={6in, 9in}]{geometry}

\usepackage[a4paper, top=1in,
bottom=1in,
left=.75in,
right=.75in]{geometry}


\begin{document}

\vspace*{3.0cm}
\begin{center}
    {\fontfamily{phv}\Huge \selectfont \textbf{MAC1140 Pre-Calculus Algebra}}
\end{center}

\vspace*{1.0cm}
\begin{center}
	{\fontfamily{phv}\Huge \selectfont \textbf{Lecture Outline}}
\end{center}

\vspace*{1.0cm}
\begin{center}
	{\fontfamily{phv}\Huge \selectfont \textbf{All Instructors}}
\end{center}
\tableofcontents
\clearpage

%\begingroup
%\pagestyle{empty}%
%\cleardoublepage
%\endgroup


\setcounter{chapter}{-1}
\pagenumbering{arabic}
\pagestyle{fancy}

\newpage


\chapter{Review}
\section{Complex Numbers (ALEKS R.6)}

One property of any given real number is that its square is always nonnegative. That is, there is no real number $x$ such that $x^2 = -1$. To remedy this situation, we introduce a new number called an imaginary unit.

\begin{ColoredDefinition}{Imaginary Unit}{def:imaginary-unit}
The {\bf imaginary unit}, $i$, is defined to be a solution to the equation $x^2 = -1$. The solution set for the equation $x^2=-1$ is $\{i,-i\}$, where $i$ is the {\bf principal square root} of $-1$.
Thus, $i^2=-1$ by definition, and we use the notation $i = \sqrt{-1}$.
\end{ColoredDefinition}

\settasks{
	after-item-skip=1in
}

\begin{example}
	Simplify the following.
\end{example}

\begin{tasks}(2)
	\task \quad $\sqrt{-12}$
	\task \quad $-\sqrt{-25}$
\end{tasks}
\vfill

\begin{objectivebox}{Simplifying imaginary numbers}{}
	The powers of $i$ repeat with every fourth power. Even powers of $i$ result in $1$ or $-1$, whereas odd powers of $i$ result in $i$ or $-i$.
\end{objectivebox}

\begin{example}
	Simplify the given powers of $i$.
\end{example}

\begin{tasks}(2)
	\task \quad $i^{3}$
	\task \quad $i^{4}$
	\task \quad $i^{5}$
	\task \quad $i^{6}$
	\task \quad $i^{103}$
	\task \quad $i^{-30}$
\end{tasks}

\vfill
\pagebreak

\begin{ColoredDefinition}{Complex Number}{def:complex-number}
	The \textbf{standard form} of a \textbf{complex number} is $a+bi$, where $a$ and $b$ are real numbers. The value $a$ is called the \textbf{real part} and the value $b$ is called the \textbf{imaginary part}.
\end{ColoredDefinition}

\begin{objectivebox}{Operations with complex numbers}{}
	Addition, subtraction, multiplication, and division of complex numbers follows a natural set of rules, as discussed below.
\end{objectivebox}

\begin{example}
	Write the following numbers in standard $a + bi$ form.
\end{example}

\settasks{
	after-item-skip=2.3in
}

\begin{tasks}(2)
	\task \quad $4i-(3+2i)+(1-3i)$
	\task \quad $(4+2i)(3-i)$
 
 \addvspace{30mm}
 
	\task \quad $(3 + 4i)^{2}$
        \task \quad $\sqrt{-3}\cdot\sqrt{-12}$
\end{tasks}

\newpage

\begin{ColoredDefinition}{Complex Conjugate}{def:complex-conjugate}
	The {\bf complex conjugate} of $z=a+bi$ is $\overline{z} =a-bi$.
	Notice that the product of conjugates is a real number.  In particular,
	\begin{center}
		$z\overline{z}=(a+bi)(a-bi)$ = \hrulefill{}.
	\end{center}
\end{ColoredDefinition}

\begin{example}
	Divide the following complex numbers and write answer in $a+bi$ form.
\end{example}

\settasks{
	after-item-skip=3.2in
}

\begin{tasks}(2)
	\task \quad $\displaystyle{4+2i\over 3-i}$
	\task \quad $\displaystyle{4\over 3+5i}$
        \task \quad $\displaystyle{\sqrt{-24}\over \sqrt{-8}}$
\end{tasks}

\newpage
\textbf{Students - Check your understanding}

\begin{checkyourunderstanding}
	Write down the following as a complex number in its standard form. Leave no negative numbers under radicals and no radicals in denominator. 
\end{checkyourunderstanding}

\settasks{
	after-item-skip=0.85in
}

\begin{tasks}(3)
	\task \quad $-\sqrt{-17}$
	\task \quad $\sqrt{-5}\cdot\sqrt{-20}$
	\task \quad $\displaystyle\frac{\sqrt{-81}}{\sqrt{3}}$
\end{tasks}
\vfill
\begin{checkyourunderstanding}
	Simplify the following powers of $i$.
\end{checkyourunderstanding}
\begin{tasks}(3)
	\task \quad $i^{48}$
	\task \quad $i^{23}$
	\task \quad $i^{-19}$
\end{tasks}
\vfill
\begin{checkyourunderstanding}
	Perform the indicated operations. Write answer in $a+bi$ form.
\end{checkyourunderstanding}

\settasks{
	after-item-skip=1in
}

\begin{tasks}(1)
	\task \quad $\left( 8-3i\right) -\left( 2+4i\right) +\left( 5+7i\right)$
	\task \quad $\left( -5+4i\right) \left( 3-i\right)$
	\task \quad $\left( 10-3i\right) \left( 10+3i\right)$
	\task \quad $\displaystyle\frac{5+6i}{2-7i}$
\end{tasks}
\vfill
\newpage


\begin{objectivebox}{Solving polynomials over the set of complex numbers}{}
The quadratic equation \(ax^2+bx+c=0\) has either one or two unique solution(s), given by \[\displaystyle x={-b\pm\sqrt{b^2-4ac}\over 2a}.\]
Solving higher order polynomials involves other techniques (see Chapter 2). For instance, the expressions representing the sum and difference of two cubes can be factored as follows:
\begin{align*}
	a^3 + b^3 & = (a + b)(a^2 - ab + b^2) \\
	a^3 - b^3 & = (a - b)(a^2 + ab + b^2)
\end{align*}

\end{objectivebox}

%# Done ___________________________________________________________________________________


\begin{example}
	Find all complex zeros of $f(x) = x^4+27x$.
\end{example}
\vskip 1.5in
\vfill

\begin{example}
	Solve $(2x+5)^{2} - 7(2x+5) - 30 = 0$
\end{example}
\vskip 1.5in
\vfill

\newpage

\begin{ColoredDefinition}{Discriminant}{def:discriminant}
	The discriminant of the quadratic equation $ax^2+bx+c=0$ is \hrulefill{}.\\ If the {\bf discriminant} of a quadratic equation is

\settasks{
	after-item-skip=0.1in
}

\begin{tasks}(1)
	\task \quad zero, there is \hrulefill{}.
	\task \quad positive, there are \hrulefill{}.
	\task \quad negative, there are \hrulefill{}.
\end{tasks}
\end{ColoredDefinition}

\begin{objectivebox}{Using the discriminant to determine nature of solutions}{}
	The discriminant is a number associated with quadratic equations, which tells us the number and type of solutions to the equation.
\end{objectivebox}

\begin{example}
	Use the discriminant to determine the number and type of solutions for each equation.
\end{example}

\settasks{
	after-item-skip=1.5in
}

\begin{tasks}(1)
	\task \quad $5x^{2}-3x+1=0$
	\task \quad $2x^{2}=3-6x$
	\task \quad $4x^{2}+12x=-9$
\end{tasks}

\vfill
\newpage

\textbf{Students - Check your understanding}

\begin{checkyourunderstanding}
	Solve the following equations.
\end{checkyourunderstanding}

\settasks{
	after-item-skip=3in
}

\begin{tasks}(1)
	\task \quad $\frac{3}{10}x^{2}-\frac{2}{5}x+\frac{7}{10}=0$
	\task \quad $x^4 - 5x^2 + 6 = 0$.
\end{tasks}
\vskip 3in

\begin{checkyourunderstanding}
	Find the discriminant and give the number and type of solution(s) of the quadratic equation $2x^2-3x+7=0$.
\end{checkyourunderstanding}

\end{document}